%===============================================================%
%  MD. RAYHAN HOSSAIN — two-column CV                           %
%  Compile with pdfLaTeX. All packages are standard on Overleaf.%
%===============================================================%
\documentclass[a4paper,10pt]{article}

\usepackage[T1]{fontenc}
\usepackage[utf8]{inputenc}
\usepackage[left=0.5in,right=0.5in,top=0.45in,bottom=0.45in]{geometry}
\usepackage{XCharter}
\usepackage{xcolor}
\usepackage{paracol}
\usepackage{enumitem}
\usepackage[hidelinks]{hyperref}

\definecolor{accent}{HTML}{18697A}
\definecolor{ink}{HTML}{1A1A1A}
\definecolor{soft}{HTML}{555F63}
\definecolor{hairline}{HTML}{BFCDD1}

\pagestyle{empty}
\raggedright
\setlength{\parindent}{0pt}
\setlength{\parskip}{0pt}
\color{ink}

% ---- separator -------------------------------------------------
\newcommand{\dotsep}{\ \textperiodcentered\ }

% ---- section heading -------------------------------------------
\newcommand{\cvsection}[1]{%
  \vspace{7pt}%
  {\color{accent}\bfseries\scshape\large #1}\par
  \vspace{1pt}%
  \noindent\textcolor{hairline}{\rule{\linewidth}{0.7pt}}\par
  \vspace{4pt}}

% ---- experience / project entry --------------------------------
% #1 title   #2 organisation or role   #3 dates   #4 stack line
\newcommand{\cventry}[4]{%
  {\bfseries #1}\par
  \vspace{0.5pt}{\color{accent}\small #2}\hfill{\small\itshape #3}\par
  \vspace{1pt}{\footnotesize\color{soft} #4}\par
  \vspace{2pt}}

% ---- skills group (right column) -------------------------------
\newcommand{\skillgroup}[2]{%
  {\bfseries\small #1}\par{\small #2}\par\vspace{4pt}}

% ---- award / certification line --------------------------------
\newcommand{\award}[2]{%
  {\small\bfseries #1}\par{\footnotesize\color{soft} #2}\par\vspace{3.5pt}}

\setlist[itemize]{leftmargin=1.05em, itemsep=1.6pt, topsep=2.5pt,
                  parsep=0pt, partopsep=0pt,
                  label=\textcolor{accent}{\small\textbullet}}

\columnratio{0.635}
\setlength{\columnsep}{7mm}

\begin{document}

\begin{center}
  {\Huge\bfseries MD. RAYHAN HOSSAIN}\par\vspace{4pt}
  {\large\color{accent}\scshape Backend Engineer \textbar\ Python, Django, FastAPI}\par\vspace{6pt}
  {\small mostafaizurrahman2021@gmail.com \dotsep 01312961737 \dotsep Dhaka, Bangladesh}\par\vspace{2.5pt}
  {\small \href{https://www.linkedin.com/in/mdrayhan03/}{linkedin.com/in/mdrayhan03} \dotsep
          \href{https://github.com/mdrayhan03}{github.com/mdrayhan03} \dotsep
          \href{https://rayhan-portfolio.lovable.app}{rayhan-portfolio.lovable.app}}
\end{center}
\vspace{4pt}\noindent\textcolor{hairline}{\rule{\textwidth}{1pt}}\par\vspace{2pt}

\begin{paracol}{2}

\cvsection{Professional Summary}
Backend Engineer specialising in Python API development at scale. Builds and optimises REST services
with Django, Django REST Framework and FastAPI, backed by PostgreSQL, Redis caching and Celery task
queues. Cut average API response time \textasciitilde35\% on a production system by restructuring
storage and asynchronous processing.

\cvsection{Experience}

\cventry{Full Stack Software Engineer}{Janatawifi LTD}{Dec 2025 -- Present}
{Python \dotsep Django \dotsep FastAPI \dotsep PostgreSQL \dotsep Redis \dotsep Celery \dotsep Docker \dotsep Azure \dotsep Frappe/ERPNext}
\begin{itemize}
  \item Cut average \textbf{API response time by \textasciitilde35\%} by offloading asset storage to
        Cloudflare R2 (S3-compatible) object storage and resolving server storage bottlenecks.
  \item Improved backend throughput using \textbf{Celery} task queues and \textbf{Redis} caching for
        \textbf{asynchronous processing} of long-running jobs.
  \item Built the \textbf{TraderBro} backend --- \textbf{FastAPI} services orchestrating LangChain and
        LangGraph agents across four LLM providers with tool calling and fallback routing for real-time
        signal aggregation.
  \item Automated the \textbf{CWTAMC} real-estate ERP operations workflow on \textbf{Frappe / ERPNext}
        (Python, MariaDB), replacing a fully manual Excel-based process and eliminating hours of
        repetitive data entry.
  \item Containerised services with \textbf{Docker} on \textbf{Azure} and monitored production
        containers (CPU, memory, storage) with structured Python logging to diagnose and resolve live
        issues.
\end{itemize}
\vspace{4pt}

\cventry{Undergraduate Teaching Assistant}{Independent University, Bangladesh}{Feb 2024 -- Nov 2025}
{Java \dotsep JavaFX \dotsep Git/GitHub}
\begin{itemize}
  \item Taught \textbf{object-oriented programming}, Java and Git/GitHub to \textbf{200+ students},
        mentoring them through design, code review and debugging on their own projects.
\end{itemize}

\cvsection{Projects}

\cventry{Plastic Kothay --- Crowdsourced Waste Reporting Platform}{Sole engineer}{May 2025 -- Present}
{Django \dotsep Django REST Framework \dotsep PostgreSQL \dotsep Docker \dotsep GitHub Actions \dotsep pytest}
{\footnotesize\color{accent}\href{https://github.com/mdrayhan03/plastickothay_dev}{github.com/mdrayhan03/plastickothay\_dev}}\par\vspace{2pt}
\begin{itemize}
  \item Architected the backend as a \textbf{hexagonal (ports-and-adapters)} design: a framework-free
        domain core, abstract ports, and swappable adapters for persistence, storage, email and
        authentication wired through a single composition root.
  \item Built a full \textbf{test pyramid} across \textbf{31 test files} --- unit tests on the pure
        domain, \textbf{integration tests} on the API, and \textbf{contract tests} asserting the domain
        scoring spec against the production SQL leaderboard query to catch drift.
  \item Enforced the architecture in \textbf{CI with GitHub Actions}: \texttt{lint-imports} blocks any
        framework import reaching the domain, with ruff and pytest on every push.
  \item Designed scoring to \textbf{derive points from current state} rather than an append-only ledger,
        so hiding a post automatically strips its points and no reconciliation job is ever required.
\end{itemize}
\vspace{4pt}

\cventry{Volunteer Portal --- Volunteer Management API}{Backend}{Nov 2025 -- Apr 2026}
{Django REST Framework \dotsep Supabase \dotsep Next.js}
\begin{itemize}
  \item Built a \textbf{REST API} managing \textbf{1000+ volunteers} across attendance tracking, profile
        management and \textbf{role-based access control}.
  \item Designed and documented \textbf{RESTful endpoints} consumed by a Next.js frontend, with DRF test
        cases covering endpoints and permission rules.
\end{itemize}

\switchcolumn

\cvsection{Education}
{\bfseries BSc, Computer Science \& Engineering}\par
{\small\color{accent} Independent University, Bangladesh}\par
{\small Sep 2022 -- \textbf{Expected Sep 2026}}\par
{\small GPA \textbf{3.81 / 4.00} \dotsep Dean's List}

\cvsection{Technical Skills}
\skillgroup{Languages}{Python, SQL, JavaScript, TypeScript, Java, Bash}
\skillgroup{Backend \& APIs}{Django, Django REST Framework (DRF), FastAPI, Flask, REST API design, Celery, JWT, OAuth 2.0, Frappe, ERPNext}
\skillgroup{Databases \& Caching}{PostgreSQL, MySQL, MariaDB, MongoDB, Redis, Qdrant, schema design, indexing, query optimization}
\skillgroup{Infrastructure}{Docker, Docker Compose, Azure, AWS, Cloudflare R2, Nginx, Linux, GitHub Actions, CI/CD}
\skillgroup{Testing}{pytest, unit testing, integration testing, contract testing, Django test framework}
\skillgroup{Architecture}{Hexagonal / clean architecture, OOP, SOLID, design patterns, system design, message queues, asynchronous processing, caching strategies}
\skillgroup{AI / LLM}{LangChain, LangGraph, OpenAI API, Anthropic API, RAG}
\skillgroup{Ways of Working}{Agile, Scrum, code review, Git/GitHub}

\cvsection{Awards}
\award{Dean's List Honoree}{Independent University Bangladesh \dotsep Dec 2022}
\award{1st Runner-Up, Hult Prize IUB}{Feb 2025}
\award{1st Runner-Up, Cisco IoT Hackathon}{Nov 2024}
\award{2nd Runner-Up, IUB Intra Programming Contest}{Nov 2023}

\cvsection{Certifications}
\award{Docker Essentials: A Developer Introduction}{IBM \dotsep Jul 2025}
\award{Agile Project Management}{HP LIFE \dotsep Jul 2025}
\award{CCNA: Introduction to Networks}{Cisco \dotsep Aug 2025}

\end{paracol}
\end{document}
